\begin{lstlisting}[language=C++,label={stan::21_05},caption={群の数を一般化したコード}]
data{
  int<lower=0> Lv;              // 水準数
  int<lower=0> N;               // 各群のサンプルサイズ       
  array[Lv,N] real X;           // 変数の値
}

parameters{
  real gm;                      // 全体平均
  array[Lv-1] real raw_delta;   // 全体からの差。水準数マイナス1個
  real<lower=0> sig;            // 誤差の分散
}

transformed parameters{
  array[Lv] real delta;           // 差の大きさを作り直す
  array[Lv] real mu;              // 再構成される群ごとの平均
  
  for(i in 1:(Lv-1)){
    delta[i] = raw_delta[i];    // ほとんどコピー
  }
  
  delta[Lv] = 0 - sum(raw_delta); // 総和が0になるように最後だけ書き換える

  for(i in 1:Lv){
    mu[i] = gm + delta[i];        // 群ごとに再構成
  }
  
}

model{
  // Likelihood
  for(l in 1:Lv){
      for(i in 1:N){
          X[l,i] ~ normal(mu[l],sig);
          
      }
  }
    
  // Prior
  gm ~ uniform(0,100);
  raw_delta ~ normal(0,100);
  sig ~ cauchy(0,5);
}

\end{lstlisting}
